% =====================================================================
%  SEMANTiCS 2026 poster  --  A0 portrait
%  Compile with:  pdflatex poster.tex
% =====================================================================
\documentclass[final]{beamer}
\usepackage[orientation=portrait,size=a0,scale=1.26]{beamerposter}
\usetheme{default}
\usefonttheme{professionalfonts}
\setbeamertemplate{navigation symbols}{}
\setbeamersize{text margin left=28mm, text margin right=28mm}
\newlength{\colh}
\AtBeginDocument{\typeout{DIMS textwidth=\the\textwidth\space textheight=\the\textheight\space paperheight=\the\paperheight}}

\usepackage[T1]{fontenc}
\usepackage[utf8]{inputenc}
\usepackage[scaled]{helvet}
\renewcommand{\familydefault}{\sfdefault}

\usepackage{booktabs}
\usepackage{tcolorbox}
\tcbuselibrary{skins,breakable}
\usepackage{graphicx}
\usepackage{ragged2e}
\usepackage{xcolor}

% ------------------------------------------------------------ palette
% Internationale huisstijl Vlaamse overheid
\definecolor{voyellow}{HTML}{FFED00}   % hoofdkleur - Pantone Yellow C
\definecolor{vodark}{HTML}{3C3D3C}     % Pantone 447
\definecolor{volight}{RGB}{246,245,243}% Cool gray 1
% aliases used throughout the body
\colorlet{navy}{vodark}
\colorlet{navydeep}{vodark}
\colorlet{teal}{vodark}
\colorlet{green}{vodark}
\colorlet{statred}{vodark}
\colorlet{orange}{voyellow}
\colorlet{panelblue}{volight}
\colorlet{panelgrey}{volight}
\colorlet{panelgold}{volight}
\colorlet{rule}{vodark}
\colorlet{ink}{vodark}

\setbeamercolor{background canvas}{bg=white}
\color{ink}

% ------------------------------------------------------------ blocks
% \begin{pbox}{<title colour>}{<body colour>}{<title>} ... \end{pbox}
\newtcolorbox{pbox}[3]{
  enhanced, breakable=false,
  colback=#2, colframe=#2, boxrule=0pt, arc=0pt, outer arc=0pt,
  left=8mm, right=8mm, top=6mm, bottom=6mm,
  title={#3}, colbacktitle=#1, coltitle=white,
  fonttitle=\bfseries\large,
  toptitle=4mm, bottomtitle=4mm, lefttitle=6mm,
  titlerule=0pt, sharp corners,
  boxsep=0pt, before skip=0pt, after skip=10mm
}

% big-number callout
\newtcolorbox{callout}[2]{
  enhanced, colback=#1, colframe=#1, boxrule=0pt, arc=0pt, sharp corners,
  left=8mm, right=8mm, top=8mm, bottom=8mm,
  borderline west={6mm}{0mm}{#2},
  before skip=0pt, after skip=10mm
}

% single statistic tile
\newcommand{\stat}[3]{%
  \begin{tcolorbox}[enhanced, colback=white, colframe=rule, boxrule=0.6mm,
      arc=0pt, sharp corners, left=2mm, right=2mm, top=4mm, bottom=4mm,
      borderline north={3mm}{0mm}{voyellow}]
    \centering{\fontsize{62}{66}\selectfont\bfseries\color{vodark}#2}\\[3mm]
    {\small #3}
  \end{tcolorbox}}

\newcommand{\figcap}[1]{\vspace{2mm}{\footnotesize\itshape #1}\vspace{2mm}}

\begin{document}
\begin{frame}[t]
\vspace*{18pt}
\setlength{\colh}{0.7116\textheight}

% ==================================================================== HEADER
\begin{tcolorbox}[enhanced, colback=voyellow, colframe=voyellow, boxrule=0pt,
    arc=0pt, sharp corners, left=14mm, right=14mm, top=10mm, bottom=10mm,
    borderline south={5mm}{0mm}{vodark}, after skip=12mm]
\begin{minipage}[c]{0.79\linewidth}
  {\fontsize{86}{96}\selectfont\bfseries\color{vodark}
   Chemical Identity Lost in Regulation}\\[5mm]
  {\fontsize{44}{50}\selectfont\itshape\color{vodark}
   A Study of Semantic Interoperability in European Chemical Substance Data}\\[9mm]
  {\fontsize{31}{36}\selectfont\bfseries\color{vodark}
   Geert Van Haute \quad Stijn Goedertier \quad Pieter Fannes \quad Maxim Van de Wynckel}\\[3mm]
  {\fontsize{25}{30}\selectfont\color{vodark}
   Department of Environment \& Spatial Development, Flemish Government, Brussels
   \enspace$\bullet$\enspace SEMANTiCS 2026, Ghent}
\end{minipage}\hfill
\begin{minipage}[c]{0.19\linewidth}\raggedleft
  
\includegraphics[width=\linewidth]{../images/State_of_the_art_horizontal_yellow_full_version_qlaz5v.jpeg}
\end{minipage}
\end{tcolorbox}

\begin{columns}[t,onlytextwidth]

% ==================================================================== COLUMN 1
\begin{column}{0.315\textwidth}
\begin{minipage}[t][\colh][s]{\linewidth}
\strut\par

\begin{pbox}{navy}{panelgrey}{The problem}
  \justifying
  In chemistry, substance identity is exact: structure-based identifiers
  (\textbf{InChI\,/\,InChIKey}) link a molecule unambiguously across PubChem,
  ChEBI and beyond.\\[6mm]
  In EU regulation (\textbf{REACH}, \textbf{CLP}), a substance may be a single
  molecule, a mixture or \textbf{UVCB}, a \textbf{group} (e.g.\ PFAS, with no
  explicit set membership), or an analytical sum parameter --- identified only
  by \textbf{CAS numbers} and textual names that carry no structure.
\end{pbox}

\includegraphics[width=0.94\linewidth]{../images/mismatch.pdf}
\figcap{Fig.~1 --- Structure-based chemical identity vs.\ heterogeneous
regulatory substance types.}

\vfill
\includegraphics[width=0.94\linewidth]{../images/regulatory_groups.pdf}
\figcap{Fig.~2 --- PFAS as a regulatory group without set-membership semantics;
Sum PFOS\,+\,PFBA a derived analytical parameter.}

\vfill
\begin{pbox}{teal}{panelblue}{Research questions}
  \textbf{RQ1}\enspace To what extent can regulatory substances be linked to a
  unique chemical structure?\\[5mm]
  \textbf{RQ2}\enspace How reliable are CAS numbers as cross-dataset
  identifiers?\\[5mm]
  \textbf{RQ3}\enspace Can embedding-based or LLM methods compensate for absent
  structure identifiers?
\end{pbox}

\begin{callout}{voyellow}{vodark}
  {\fontsize{100}{104}\selectfont\bfseries\color{vodark}37.5\%}\\[6mm]
  {\fontsize{31}{38}\selectfont\bfseries\color{vodark}
   of regulatory substance entries have no machine-readable structural identity.}
\end{callout}

\end{minipage}
\end{column}

% ==================================================================== COLUMN 2
\begin{column}{0.315\textwidth}
\begin{minipage}[t][\colh][s]{\linewidth}
\strut\par

\begin{pbox}{navy}{panelgrey}{Approach}
  \justifying
  We integrated \textbf{18 ECHA-based datasets} (41,813 records): obligation and
  activity lists, REACH registrations, the EU Pesticides Database, OSPAR and a
  Flemish administrative list. Each entry was linked to InChIKey via PubChem and
  assigned to a \textbf{seven-tier linkability taxonomy}.\\[6mm]
  Substance names were matched to ${\sim}$4,800 ChemOnt classes using four
  sentence-transformer models and Claude Sonnet 4.6, benchmarked against 28,204
  ClassyFire assignments.
\end{pbox}

\begin{minipage}[t]{0.315\linewidth}\stat{green}{62.5\%}{link to a defined structure (tier 1)}\end{minipage}\hfill
\begin{minipage}[t]{0.315\linewidth}\stat{statred}{37.5\%}{non-structure-based entries}\end{minipage}\hfill
\begin{minipage}[t]{0.315\linewidth}\stat{navy}{39.2\%}{best LLM Hit@1 (still a minority)}\end{minipage}

\vfill
\begin{pbox}{navy}{panelblue}{Linkability\ \ (RQ1)}
  \justifying
  Of 33,452 unique entries, \textbf{62.5\% (20,909)} link directly via InChIKey.
  The other \textbf{37.5\%} resist structure-based linking: 11.6\% unresolved CAS
  (tier 4), 10.8\% mixtures or UVCBs (tier 5), and 14.9\% lack any usable
  identifier (tier 7).
\end{pbox}

\includegraphics[width=0.94\linewidth]{../images/Analysis_4bd_v_Entity_type_and_linkability_vertical.pdf}
\figcap{Fig.~3 --- Entity-type composition (A) and seven-tier linkability (B)
across 18 sources.}

\vfill
\begin{pbox}{navy}{panelgrey}{CAS unreliability\ \ (RQ2)}
  \justifying
  \textbf{70.0\%} of structure-resolved InChIKeys appear in only one source ---
  CAS is insufficient as a shared key.\\[6mm]
  Ambiguity compounds this: 195 InChIKeys map to multiple CAS numbers (up to
  seven); 4 CAS numbers resolve to two distinct InChIKeys (e.g.\ toluene
  diisocyanate $\rightarrow$ two positional isomers).
\end{pbox}

\begin{callout}{vodark}{voyellow}
  {\fontsize{42}{52}\selectfont\bfseries\color{white}
   CAS numbers introduce ambiguity rather than resolving it.}
\end{callout}

\end{minipage}
\end{column}

% ==================================================================== COLUMN 3
\begin{column}{0.315\textwidth}
\begin{minipage}[t][\colh][s]{\linewidth}
\strut\par

\begin{pbox}{statred}{panelgrey}{Semantic methods fail\ \ (RQ3)}
  \justifying
  Embedding-based matching collapses: the best model reaches
  \textbf{Hit@1 = 5.75\%}; all others score below 1\%. Claude Sonnet 4.6 reaches
  \textbf{39.2\%} --- a seven-fold gain, but it still misclassifies the majority.
  Text cannot substitute for structure.
\end{pbox}

{\centering
\begin{tabular}{@{}l r r r r@{}}
\toprule
\textbf{Model} & \textbf{Hit@1} & \textbf{Hit@3} & \textbf{Hit@5} & \textbf{Hit@10}\\
\midrule
all-MiniLM-L6-v2  & 0.0575 & 0.0986 & 0.1234 & 0.1621\\
all-mpnet-base-v2 & 0.0093 & 0.0152 & 0.0184 & 0.0247\\
SciBERT           & 0.0040 & 0.0071 & 0.0102 & 0.0131\\
BioBERT           & 0.0047 & 0.0077 & 0.0099 & 0.0150\\
\textbf{Claude Sonnet 4.6} & \textbf{0.392} & --- & --- & ---\\
\bottomrule
\end{tabular}\par}
\figcap{Table~1 --- ChemOnt classification (direct-parent level). Only Hit@1 is
comparable across rows.}

\vfill
\begin{pbox}{green}{panelblue}{What works: the knowledge graph\ \ (C4, C5)}
  \justifying
  ${\sim}$17,500 regulatory substances (SKOS concepts) are integrated with
  structure-based identity, ChemOnt classification and ChEBI biological roles ---
  \textbf{${\sim}$46,600 cross-domain links}, enabling composite priority scoring
  that is infeasible on fragmented data.\\[6mm]
  LLM-based scope mapping yields \textbf{304 validated SKOS triples}
  (285 broadMatch, 18 exactMatch, 1 narrowMatch) as a legislative--ChemOnt
  crosswalk.
\end{pbox}

\includegraphics[width=0.94\linewidth]{../images/knowledge_graph.pdf}
\figcap{Fig.~4 --- Integrated knowledge graph linking identity, classification
and biological roles.}

\vfill
\begin{pbox}{navydeep}{panelgold}{Conclusions \& recommendations}
  \justifying
  Over one-third of entries cannot be linked to a structure; CAS adds ambiguity;
  neither embeddings nor LLMs can replace structural information. Manual
  reconciliation (${\sim}$3,150 person-years for REACH) is infeasible.
  \textbf{Harmonisation must happen at the source, not downstream.}\\[6mm]
  \textbf{1}\enspace Require structure-based identifiers (InChIKey; MInChI for
  mixtures) across regulatory instruments.\\[4mm]
  \textbf{2}\enspace Mandate explicit substance-group representation using
  formal ontologies such as ChemOnt.\\[4mm]
  \textbf{3}\enspace Publish regulatory data as linked data with persistent URIs
  and SKOS links (FAIR / Interoperable Europe Act).
\end{pbox}

\end{minipage}
\end{column}
\end{columns}

% ==================================================================== FOOTER
\vfill
\begin{tcolorbox}[enhanced, colback=vodark, colframe=vodark, boxrule=0pt,
    arc=0pt, sharp corners, left=14mm, right=14mm, top=7mm, bottom=7mm,
    borderline west={6mm}{0mm}{voyellow}]
\begin{minipage}[c]{0.76\linewidth}
  {\color{white}\large
   \textbf{Code \& data:}\enspace
   github.com/gezever/A-Substance-Is-Not-Always-a-Substance \\[3mm]
   \textbf{Contact:}\enspace geert.vanhaute@vlaanderen.be \\[3mm]
   {\normalsize\itshape AI declaration: Claude Sonnet 4.6 was used as the object
   of study for the LLM experiments (${\sim}$\$100 API) and as a development
   tool. All content was reviewed by the authors. CC BY 4.0.}}
\end{minipage}\hfill
\begin{minipage}[c]{0.22\linewidth}\raggedleft
  
\includegraphics[width=32mm]{../images/qrcode_github.png}\\[2mm]
  {\color{white}\normalsize Scan for code \& data}
\end{minipage}
\end{tcolorbox}
\vspace*{45pt}

\end{frame}
\end{document}
